% GENERATED FILE. Edit canonical lessons, then run npm run generate:books.
% canonical-source-hash: fnv1a64:fc534a09cc110dae
% canonical-lessons: IT-C407-parabrezza, IT-C407-tergicristallo, IT-C407-specchietto, IT-C407-clacson, IT-C407-appannarsi

\chapter{Vetri e specchietti}
\label{ch:it-vetri-e-specchietti}
\hlchaptermodality{407}

\begin{chapteropening}
\textbf{By the end of this chapter:} \emph{I can talk about seeing clearly from a car: the windscreen, the wipers, the mirror, the horn and windows that mist up.}

Complete the last lesson of chapter 407: I can talk about seeing clearly from a car: the windscreen, the wipers, the mirror, the horn and windows that mist up.
\end{chapteropening}

\section[il parabrezza]{il parabrezza — the windscreen}

\label{lesson:IT-C407-parabrezza}



\begin{quote}
Before the new one: say the Italian for to steer, then the Italian for to skid.
\end{quote}

\subsection*{You'll want to know: il parabrezza}
\textbf{il parabrezza} — \textquotedblleft{}the windscreen\textquotedblright{}.

\begin{quote}
Un sasso ha rotto il parabrezza.
\end{quote}

\begin{quote}
\emph{A stone cracked the windscreen.}
\end{quote}

Built like \emph{parasole}: \emph{para-} means it shields you from something. Plural unchanged.

\begin{cousinweb}
\emph{parabrezza} is built from \emph{parare} 'to ward off', which is inherited from Latin \emph{parāre} 'to prepare' (Lewis \& Short), and \emph{brezza} 'breeze': the glass that holds back the wind.
\end{cousinweb}

\begin{grammarlens}[title={What you've built}]
One more word for telling, explaining and describing at B1.
\end{grammarlens}

\subsection*{Your turn}
\begin{itemize}
\raggedright
  \item \emph{Say it:} \emph{il parabrezza}
  \item \emph{Say it:} \emph{il parabrezza}, once more
  \item \emph{Say it:} a sentence of your own with \emph{il parabrezza}
  \item \emph{Recall:} say the Italian for to steer, then the Italian for to skid, then say \emph{il parabrezza} again
\end{itemize}

\begin{tcolorbox}[breakable,colback=teal!4,colframe=teal!35!black,title={Before you move on}]
What does il parabrezza mean? (\textquotedblleft{}The windscreen\textquotedblright{}.) Say it once more.
\end{tcolorbox}
\section[il tergicristallo]{il tergicristallo — the windscreen wiper}

\label{lesson:IT-C407-tergicristallo}



\begin{quote}
Before the new one: say the Italian for to skid, then the Italian for the windscreen.
\end{quote}

\subsection*{You'll want to know: il tergicristallo}
\textbf{il tergicristallo} — \textquotedblleft{}the windscreen wiper\textquotedblright{}.

\begin{quote}
Accendi il tergicristallo, piove forte.
\end{quote}

\begin{quote}
\emph{Turn on the wipers, it is raining hard.}
\end{quote}

Often in the plural: \emph{i tergicristalli}.

\begin{cousinweb}
\emph{tergicristallo} is built from \emph{tergere} 'to wipe', a learned borrowing from Latin \emph{tergēre} (Lewis \& Short), and \emph{cristallo} 'glass', inherited from Latin \emph{crystallum}, from Greek \emph{krýstallos} 'clear ice' (EtymDB).
\end{cousinweb}

\begin{grammarlens}[title={What you've built}]
One more word for telling, explaining and describing at B1.
\end{grammarlens}

\subsection*{Your turn}
\begin{itemize}
\raggedright
  \item \emph{Say it:} \emph{il tergicristallo}
  \item \emph{Say it:} \emph{il tergicristallo}, once more
  \item \emph{Say it:} a sentence of your own with \emph{il tergicristallo}
  \item \emph{Recall:} say the Italian for to skid, then the Italian for the windscreen, then say \emph{il tergicristallo} again
\end{itemize}

\begin{tcolorbox}[breakable,colback=teal!4,colframe=teal!35!black,title={Before you move on}]
What does il tergicristallo mean? (\textquotedblleft{}The windscreen wiper\textquotedblright{}.) Say it once more.
\end{tcolorbox}
\section[lo specchietto]{lo specchietto — the wing mirror, the car mirror}

\label{lesson:IT-C407-specchietto}



\begin{quote}
Before the new one: say the Italian for the windscreen, then the Italian for the windscreen wiper.
\end{quote}

\subsection*{You'll want to know: lo specchietto}
\textbf{lo specchietto} — \textquotedblleft{}the wing mirror, the car mirror\textquotedblright{}.

\begin{quote}
Guarda nello specchietto prima di girare.
\end{quote}

\begin{quote}
\emph{Check the mirror before you turn.}
\end{quote}

A little \emph{specchio}; \emph{lo specchietto retrovisore} is the rear-view mirror.

\begin{cousinweb}
\emph{specchietto} is a little \emph{specchio} 'mirror', which is inherited from Vulgar Latin \emph{speclum}, from Latin \emph{speculum} 'mirror' (EtymDB).
\end{cousinweb}

\begin{grammarlens}[title={What you've built}]
One more word for telling, explaining and describing at B1.
\end{grammarlens}

\subsection*{Your turn}
\begin{itemize}
\raggedright
  \item \emph{Say it:} \emph{lo specchietto}
  \item \emph{Say it:} \emph{lo specchietto}, once more
  \item \emph{Say it:} a sentence of your own with \emph{lo specchietto}
  \item \emph{Recall:} say the Italian for the windscreen, then the Italian for the windscreen wiper, then say \emph{lo specchietto} again
\end{itemize}

\begin{tcolorbox}[breakable,colback=teal!4,colframe=teal!35!black,title={Before you move on}]
What does lo specchietto mean? (\textquotedblleft{}The wing mirror, the car mirror\textquotedblright{}.) Say it once more.
\end{tcolorbox}
\section[il clacson]{il clacson — the horn (of a car)}

\label{lesson:IT-C407-clacson}



\begin{quote}
Before the new one: say the Italian for the windscreen wiper, then the Italian for the wing mirror.
\end{quote}

\subsection*{You'll want to know: il clacson}
\textbf{il clacson} — \textquotedblleft{}the horn (of a car)\textquotedblright{}.

\begin{quote}
Il taxi suona il clacson per strada.
\end{quote}

\begin{quote}
\emph{The taxi honks its horn in the street.}
\end{quote}

\emph{Suonare il clacson} — to honk. Plural unchanged.

\begin{cousinweb}
\emph{clacson} was borrowed from English \emph{klaxon} (EtymDB), first the trade name of an electric car horn, spelled the Italian way.
\end{cousinweb}

\begin{grammarlens}[title={What you've built}]
One more word for telling, explaining and describing at B1.
\end{grammarlens}

\subsection*{Your turn}
\begin{itemize}
\raggedright
  \item \emph{Say it:} \emph{il clacson}
  \item \emph{Say it:} \emph{il clacson}, once more
  \item \emph{Say it:} a sentence of your own with \emph{il clacson}
  \item \emph{Recall:} say the Italian for the windscreen wiper, then the Italian for the wing mirror, then say \emph{il clacson} again
\end{itemize}

\begin{tcolorbox}[breakable,colback=teal!4,colframe=teal!35!black,title={Before you move on}]
What does il clacson mean? (\textquotedblleft{}The horn (of a car)\textquotedblright{}.) Say it once more.
\end{tcolorbox}
\section[appannarsi]{appannarsi — to mist up, to fog up}

\label{lesson:IT-C407-appannarsi}



\begin{quote}
Before the new one: say the Italian for the wing mirror, then the Italian for the horn.
\end{quote}

\subsection*{You'll want to know: appannarsi}
\textbf{appannarsi} — \textquotedblleft{}to mist up, to fog up\textquotedblright{}.

\begin{quote}
Con il freddo i vetri si appannano.
\end{quote}

\begin{quote}
\emph{In the cold the windows mist up.}
\end{quote}

\emph{Appannato} — misty; eyes and memories can be \emph{appannati} too.

\begin{cousinweb}
\emph{appannarsi} is built on \emph{panno} 'cloth', which is inherited from Latin \emph{pannus} 'cloth, rag' (EtymDB): glass that looks covered by a veil of cloth.
\end{cousinweb}

\begin{grammarlens}[title={What you've built}]
One more word for telling, explaining and describing at B1.
\end{grammarlens}

\subsection*{Your turn}
\begin{itemize}
\raggedright
  \item \emph{Say it:} \emph{appannarsi}
  \item \emph{Say it:} \emph{appannarsi}, once more
  \item \emph{Say it:} a sentence of your own with \emph{appannarsi}
  \item \emph{Recall:} say the Italian for the wing mirror, then the Italian for the horn, then say \emph{appannarsi} again
\end{itemize}

\begin{tcolorbox}[breakable,colback=teal!4,colframe=teal!35!black,title={Before you move on}]
What does appannarsi mean? (\textquotedblleft{}To mist up, to fog up\textquotedblright{}.) Say it once more.
\end{tcolorbox}
